\documentclass[10pt]{article}
\usepackage[margin=1.6cm]{geometry}
\usepackage{graphicx}
\usepackage{amsmath}
\usepackage{amssymb}
\usepackage[hidelinks]{hyperref}
\setlength{\parskip}{2pt}\setlength{\parindent}{0pt}

\title{\vspace{-1.4cm}Block-2--only weight decay reverses the story:\\
protected parents, rising $R(t)$, and decay that again favors composition\\
\large exp.~6 --- identical to exp.~4 except phase-2 $\lambda$ acts on $V_2$ only ---
8 $\alpha$ $\times$ 8 $\lambda$ $\times$ 3 seeds}
\author{}\date{\vspace{-1.2cm}}

\begin{document}\maketitle\vspace{-1.9cm}

\section*{Design}
Identical to exp.~4 (converged-$d2$ pretrain, $V_1\!\to\!\alpha V_1$ at $d4$ onset setting
$R_0\in[0.04,0.57]$, everything trainable, $d2$ losses kept, fresh optimizer, exact Walsh,
per-cell measurements every 500 steps) with \textbf{one change}: the phase-2 weight decay is
applied \emph{only to the second-block weights $V_2$}; embedding, $V_1$ and readout are
decay-free. The circuit under study is thus regularized, while the representation that feeds it
is protected --- the control that exp.~4 identified as missing. The 10k pretraining checkpoints
are saved per seed (\texttt{accessibility\_ckpts/}, also on Drive). Every cell solves $d4$
($|\hat f_4|=1.000$ throughout).

\section*{Findings}
\textbf{1. Decay again favors the compositional circuit, with the boundary at $R_c$.} Right of
$R_c$ the parent share now \emph{grows} with $\lambda$: at $R_0=0.52$,
$f_{2+2}=0.103\;(\lambda{=}0)\to0.175\;(10^{-3})\to0.230\;(10^{-2})\to0.385\;(3{\cdot}10^{-2})$;
left of $R_c$ it stays $\leq0.01$ at every $\lambda$ (Figs.~\ref{fig:phase},~\ref{fig:drift}
left). This is the frozen-case direction restored --- global decay (exp.~4) drove the same
cells to $f_{2+2}=0$. The magnitude still trails the frozen benchmark ($0.385$ vs.\ $0.90$ at
matched $R,\lambda$): with the task solved within the first few hundred steps, the early
singleton lock-in owns most of the output, and $V_2$-decay reworks only part of it.

\textbf{2. The representation drift reverses.} With $V_1$ protected, $\lambda$ now pushes
$R_{\rm final}$ \emph{up}, not down (Fig.~\ref{fig:drift}, right): from $R_0=0.04$,
$R_{\rm final}$ rises $0.14\to0.29$ as $\lambda\!:0\to10^{-2}$; supercritical starts stay
supercritical at every $\lambda$ ($R_{\rm final}\approx0.41$--$0.46$ at $R_0=0.52$, vs.\
$0.21$ under global decay). Mechanically: taxing block-2 circuits shifts gradient pressure
into the untaxed $V_1$/embedding, strengthening the parents --- the opposite of exp.~4's
inverse-$\alpha$ erosion.

\textbf{3. No dismantling transients.} The phase plane (Fig.~\ref{fig:plane}) shows
trajectories moving \emph{up} in $f_{2+2}$ under decay (ends above starts), never the
build-then-demolish loops of exp.~4: with $R(t)$ held supercritical, decay works \emph{for}
the compositional circuit throughout training.

\textbf{4. Strong decay exits through a different door.} At $\lambda=10^{-1}$ (all cells) and
a few at $3{\cdot}10^{-2}$, block~2 is suppressed entirely --- no neuron contributes to $d4$
($n_{\rm rel}=0$; NaN cells, dropped from the diagrams) --- yet $\hat f_4=1.0$: the readout
solves the task directly from the frozen stream's own $d4$ leak
($\|q_{1234}\|\approx0.105\,\alpha$ from the block-1 corner units), the same bypass seen at
$\alpha=8$ in the frozen experiment. The ``circuit choice'' question dissolves there rather
than being answered.

\begin{figure}[h!]\centering
\includegraphics[width=0.49\textwidth]{v2_phase_R0.pdf}\hfill
\includegraphics[width=0.49\textwidth]{v2_phase_Rfinal.pdf}\vspace{-4pt}
\caption{Phase diagrams vs.\ $R_0$ (left) and $R_{\rm final}$ (right), colors on the same
scale as exps.~2/4. Right of $R_c$ the color now deepens toward composition as $\lambda$
grows (up the page); missing squares in the top rows are the block-2--bypassed (NaN) cells.}
\label{fig:phase}
\end{figure}

\begin{figure}[h!]\centering
\includegraphics[width=0.49\textwidth]{v2_cuts_fixed_lambda.pdf}\hfill
\includegraphics[width=0.49\textwidth]{v2_drift.pdf}\vspace{-4pt}
\caption{Left: $f_{2+2}$ vs.\ $R_0$ at fixed $\lambda$ --- ordered by $\lambda$ right of
$R_c$ (compare exp.~4, where the order was reversed and collapsed to 0). Right: drift map ---
$\lambda$ raises $R_{\rm final}$ everywhere; supercritical stays supercritical.}
\label{fig:drift}
\end{figure}

\begin{figure}[h!]\centering
\includegraphics[width=0.98\textwidth]{v2_phase_plane.pdf}\vspace{-4pt}
\caption{Phase-plane trajectories (seed 0; $\circ$ start, $\blacksquare$ end). Unlike exp.~4
there are no dismantling loops: under $V_2$-decay, $f_{2+2}$ grows along the trajectory and
$R(t)$ never collapses.}
\label{fig:plane}
\end{figure}

\textbf{Bottom line.} Exp.~4's conclusion sharpens: unfrozen dynamics per se do not kill the
compositional phase --- \emph{decay on the representation} does. Protecting $V_1$ restores
weight decay's frozen-case role qualitatively (sharpening toward the efficient circuit, pinned
at $R_c$, now with the representation itself drifting upward in accessibility). The remaining
factor-$\sim$2 gap to the frozen benchmark is purely dynamical --- the direct circuit forms
first and is only partially reworked --- so closing it should require slowing $d4$ learning
(smaller readout lr) or annealed decay, not more decay.

\vspace{2pt}\footnotesize
Artifacts: \texttt{gdrive:parity/analysis\_output/accessibility\_alpha\_v2wd/}
(\texttt{results\_merged.csv}, 12 worker slices, trajectories, figures, this report);
checkpoints: \texttt{accessibility\_ckpts/pretrain10k\_seed\{0,1,2\}.pt}; code:
\texttt{analysis/accessibility/run\_alpha\_unfrozen\_experiment.py} (\texttt{--wd-block2-only}),
\texttt{merge\_alpha\_unfrozen.py}, \texttt{scripts/launch\_v2wd\_workers.sh}. Companions:
exps.~1--4 reports.
\end{document}
